%! Polynomial Long and Synthetic Division — Unit 3, Lesson 3.5 — Slides (Beamer)
% Release order: title, hook, I Do (rows LONG DIVISION, SYNTHETIC), We Do (REMAINDER THEOREM,
% FACTOR THEOREM and the missing-power crux), You Do (DIVIDE & FACTOR) -> debrief, closing on
% the DeltaMath assignment.
\documentclass[aspectratio=169, 11pt]{beamer}
\usepackage{precalculus-beamer}   % provides \plumheader{} and \sectionlabel[color]{}

\begin{document}

% TITLE SLIDE
{
\setbeamercolor{background canvas}{bg=plum}
\begin{frame}[plain]
  \vspace{2.6cm}\hspace{0.55cm}%
  \begin{minipage}{0.9\textwidth}
    {\color{goldacc}\bfseries\small LESSON 3.5}\\[0.5cm]
    {\color{white}\bfseries\Large Polynomial Long and}\\[0.2cm]
    {\color{white}\bfseries\Large Synthetic Division}\\[1.0cm]
    {\color{lilacmid}\small Precalculus~$\cdot$~Unit 3: Polynomial Functions}
  \end{minipage}
\end{frame}
}

% HOOK
\begin{frame}[t, plain]
  \plumheader{Hook: No Pattern Fits}
  \sectionlabel[goldacc]{Think About It}

  \begin{center}
    {\Large $p(x) = x^3 - 7x + 6$}
  \end{center}

  No common factor. No grouping. Not a sum or difference of cubes.

  \bigskip
  \begin{columns}[T]
    \column{0.48\textwidth}
      \begin{block}{Maya}
        ``Then there is no way to tell whether $x - 2$ is a factor.''
      \end{block}
    \column{0.48\textwidth}
      \begin{block}{Jordan}
        ``Just plug in 2.''
      \end{block}
  \end{columns}

  \bigskip
  \begin{center}
    {\large\bfseries\color{plum} How could plugging in one number tell you anything about a factor?}
  \end{center}
\end{frame}

% ROW 1 — LONG DIVISION
\begin{frame}[t, plain]
  \plumheader{Long Division --- Like Dividing Numbers}
  \sectionlabel[goldacc]{Notes row LONG DIVISION $\cdot$ I Do: problem 1 $\cdot$ We Do: problem 2}

  \begin{columns}[T]
    \column{0.30\textwidth}
      \begin{center}
      {\setlength{\tabcolsep}{3pt}\begin{tabular}{r@{\,}ccc}
                   &        & \textbf{1} & \textbf{3} \\ \cline{2-4}
        $12\,\big)$ & $1$    & $5$    & $7$    \\
                   & $-(1$  & $2)$   &        \\ \cline{2-3}
                   &        & $3$    & $7$    \\
                   &        & $-(3$  & $6)$   \\ \cline{3-4}
                   &        &        & \textbf{\color{redacc}1} \\
      \end{tabular}}

      \medskip
      $157 = 12 \cdot 13 + 1$
      \end{center}
    \column{0.36\textwidth}
      \begin{center}
      {\setlength{\tabcolsep}{3pt}\begin{tabular}{r@{\,}ccc}
                     &          & $\mathbf{x}$ & $\mathbf{+3}$ \\ \cline{2-4}
        $x+2\,\big)$ & $x^2$    & $+5x$    & $+7$    \\
                     & $-(x^2$  & $+2x)$   &         \\ \cline{2-3}
                     &          & $3x$     & $+7$    \\
                     &          & $-(3x$   & $+6)$   \\ \cline{3-4}
                     &          &          & {\color{redacc}$\mathbf{1}$} \\
      \end{tabular}}

      \medskip
      $x^2 + 5x + 7 = (x + 2)(x + 3) + 1$
      \end{center}
    \column{0.30\textwidth}
      \begin{block}{Division algorithm}
        $f(x) = g(x)\,q(x) + r(x)$\\[3pt]
        remainder's degree $<$ divisor's degree
      \end{block}
      \smallskip
      {\small Divide $\to$ multiply $\to$ subtract (change every sign) $\to$ bring down.}
  \end{columns}
\end{frame}

% ROW 2 — SYNTHETIC
\begin{frame}[t, plain]
  \plumheader{Synthetic Division --- The Same Numbers, Faster}
  \sectionlabel[goldacc]{Notes row SYNTHETIC $\cdot$ I Do: problem 3 $\cdot$ We Do: problems 4--6}

  \begin{columns}[T]
    \column{0.46\textwidth}
      \begin{center}
      {\setlength{\tabcolsep}{2.5pt}\begin{tabular}{r@{\,}cccc}
                   &          & $\mathbf{x^2}$ & $\mathbf{+x}$ & $\mathbf{-2}$ \\ \cline{2-5}
       $x-3\,\big)$ & $x^3$    & $-2x^2$  & $-5x$    & $+6$   \\
                   & $-(x^3$  & $-3x^2)$ &          &        \\ \cline{2-3}
                   &          & $x^2$    & $-5x$    &        \\
                   &          & $-(x^2$  & $-3x)$   &        \\ \cline{3-4}
                   &          &          & $-2x$    & $+6$   \\
                   &          &          & $-(-2x$  & $+6)$  \\ \cline{4-5}
                   &          &          &          & {\color{redacc}$\mathbf{0}$} \\
      \end{tabular}}
      \end{center}
    \column{0.50\textwidth}
      \begin{center}
      {\large\setlength{\tabcolsep}{6pt}\begin{tabular}{c|cccc}
        $3$ & $1$ & $-2$ & $-5$ & $6$ \\
            &     & $3$  & $3$  & $-6$ \\ \hline
            & $\mathbf{1}$ & $\mathbf{1}$ & $\mathbf{-2}$ & {\color{redacc}$\mathbf{0}$} \\
      \end{tabular}}
      \end{center}

      \medskip
      \begin{itemize}\setlength{\itemsep}{4pt}
        \item Only for a divisor $x - c$. \quad $x + 2 \to c = -2$.
        \item Bring down; multiply by $c$; add; repeat.
        \item Last number: the \textbf{remainder}. The rest: the quotient's
              \textbf{coefficients}, one degree lower.
      \end{itemize}
  \end{columns}
\end{frame}

% ROW 3 — REMAINDER THEOREM
\begin{frame}[t, plain]
  \plumheader{The Remainder Theorem --- Two Roads to $p(c)$}
  \sectionlabel[goldacc]{Notes row REMAINDER THEOREM $\cdot$ I Do: problem 7 $\cdot$ We Do: problem 8}

  \begin{block}{Remainder Theorem}
    When $p(x)$ is divided by $x - c$, the remainder is $p(c)$.
  \end{block}

  \medskip
  \begin{columns}[T]
    \column{0.44\textwidth}
      \textbf{Road 1: substitute.} \quad $p(x) = x^3 - 2x^2 + 4$
      \[ p(3) = 27 - 18 + 4 = \mathbf{13} \]
    \column{0.48\textwidth}
      \textbf{Road 2: divide by $x - 3$.}

      \smallskip
      \begin{center}
      {\setlength{\tabcolsep}{5pt}\begin{tabular}{c|cccc}
        $3$ & $1$ & $-2$ & $0$ & $4$ \\
            &     & $3$  & $3$ & $9$ \\ \hline
            & $1$ & $1$  & $3$ & {\color{redacc}$\mathbf{13}$} \\
      \end{tabular}}
      \end{center}
  \end{columns}

  \medskip
  \textbf{Why:} $p(x) = (x - 3)\,q(x) + r$, so $p(3) = 0 \cdot q(3) + r = r$.
\end{frame}

% ROW 4 — FACTOR THEOREM
\begin{frame}[t, plain]
  \plumheader{The Factor Theorem --- Remainder Zero}
  \sectionlabel[goldacc]{Notes row FACTOR THEOREM $\cdot$ I Do: problem 9 $\cdot$ We Do: problem 10}

  \begin{columns}[T]
    \column{0.46\textwidth}
      \begin{center}
      \begin{tikzpicture}[x=0.6cm, y=0.1cm]
        \draw[linegray, very thin, xstep=1, ystep=5] (-4,-15) grid (3,15);
        \draw[->, thick] (-4,0) -- (3.5,0) node[right] {\small $x$};
        \draw[->, thick] (0,-15) -- (0,16.5) node[above] {\small $p(x)$};
        \foreach \x in {-3,-2,-1,1,2,3} \node[below, font=\scriptsize] at (\x,0) {\x};
        \draw[very thick, plum, domain=-3.55:2.95, samples=80, smooth]
             plot (\x, {\x*\x*\x - 7*\x + 6});
        \foreach \x in {-3,1,2} \fill[plum] (\x,0) circle (2.5pt);
      \end{tikzpicture}

      {\footnotesize $p(x) = x^3 - 7x + 6$}
      \end{center}
    \column{0.50\textwidth}
      \begin{block}{Factor Theorem}
        $(x - c)$ is a factor of $p(x)$ exactly when $p(c) = 0$.
      \end{block}
      \medskip
      remainder $0$ $\iff$ $(x - c)$ is a factor\\[2pt]
      \hspace*{3.2em}$\iff$ $c$ is a zero $\iff$ $(c, 0)$ is an $x$-intercept

      \medskip
      \textbf{One zero known?} Divide its factor out, then factor the quadratic quotient.
  \end{columns}
\end{frame}

% THE CRUX
\begin{frame}[t, plain]
  \plumheader{Is $x - 2$ a Factor of $x^3 - 7x + 6$?}
  \sectionlabel[redacc]{Notes problems 11--14 $\cdot$ two divisions that disagree}

  \begin{columns}[T]
    \column{0.46\textwidth}
      \begin{block}{Ana}
        \begin{center}
        {\large\setlength{\tabcolsep}{6pt}\begin{tabular}{c|ccc}
          $2$ & $1$ & $-7$ & $6$ \\
              &     & $?$  & $?$ \\ \hline
              & $?$ & $?$  & $?$ \\
        \end{tabular}}
        \end{center}
      \end{block}
    \column{0.46\textwidth}
      \begin{block}{Ben}
        \begin{center}
        {\large\setlength{\tabcolsep}{6pt}\begin{tabular}{c|cccc}
          $2$ & $1$ & $0$ & $-7$ & $6$ \\
              &     & $?$ & $?$  & $?$ \\ \hline
              & $?$ & $?$ & $?$  & $?$ \\
        \end{tabular}}
        \end{center}
      \end{block}
  \end{columns}

  \bigskip
  \begin{itemize}\setlength{\itemsep}{5pt}
    \item Finish both. Do the remainders agree?
    \item Settle it with $p(2)$. Whose remainder does the Remainder Theorem back?
    \item Multiply out Ana's answer. Which polynomial did she \emph{really} divide?
  \end{itemize}

  \medskip
  \textbf{A missing power still gets a column: write $0$ for it.}
\end{frame}

% YOU DO
\begin{frame}[t, plain]
  \plumheader{You Do --- Divide \& Factor}
  \sectionlabel[redacc]{Notes row DIVIDE \& FACTOR $\cdot$ problems 15--18, alone, 17 minutes}

  \begin{columns}[T]
    \column{0.46\textwidth}
      \begin{block}{Checklist, every time}
        \textbf{1.} Every power, highest first --- $0$ for a missing one.\\[3pt]
        \textbf{2.} $x - c$: use $c$. \ $x + c$: use $-c$.\\[3pt]
        \textbf{3.} Last number $=$ remainder $= p(c)$.
      \end{block}
    \column{0.50\textwidth}
      \begin{enumerate}\setlength{\itemsep}{6pt}
        \setcounter{enumi}{14}
        \item $(2x^3 + 3x^2 - 4x + 5) \div (x + 2)$
        \item Remainder of $x^4 - 3x^2 + 5$ divided by $x - 2$, two ways.
        \item Show $x - 4$ is a factor of $x^3 - 13x - 12$; find all zeros.
        \item Find $k$ so that $x + 2$ is a factor of $x^3 + kx^2 - x - 6$.
      \end{enumerate}
  \end{columns}
\end{frame}

% DEBRIEF
\begin{frame}[t, plain]
  \plumheader{Debrief --- The Headline}
  \sectionlabel[goldacc]{Share out 15--18, one answer each $\cdot$ then say it back}

  \begin{columns}[T]
    \column{0.48\textwidth}
      \begin{block}{Remainder and factor}
        Dividing by $x - c$ leaves the remainder $p(c)$.\\[3pt]
        Remainder $0$ $\Rightarrow$ $x - c$ is a factor and $c$ is a zero.
      \end{block}
    \column{0.48\textwidth}
      \begin{block}{Every column counts}
        A missing power gets a $0$.\\[3pt]
        Without it, you divide a different polynomial.
      \end{block}
  \end{columns}

  \bigskip
  \textbf{Next time (3.6):} the chain runs backwards --- start from the zeros, write the factors,
  and multiply out to the polynomial.
\end{frame}

% HOMEWORK — DELTAMATH
\begin{frame}[t, plain]
  \plumheader{Homework: DeltaMath}
  \sectionlabel[redacc]{Graded $\cdot$ assignment \textit{Lesson 3.5}}

  \begin{itemize}\setlength{\itemsep}{8pt}
    \item On DeltaMath, open the assignment \textbf{Lesson 3.5}. Copy the due date into the
          \textbf{Homework} row of your cover sheet: \quad Due: \rule{3cm}{0.4pt}
    \item Topics: polynomial long division $\cdot$ synthetic division $\cdot$ the Remainder
          Theorem $\cdot$ the Factor Theorem $\cdot$ synthetic division to find zeros. Before
          every division: is every power there?
    \item \textbf{AP Practice} (in the packet) is \textbf{extra credit, scored 1 to 5}. Optional.
          Do DeltaMath first.
  \end{itemize}
\end{frame}

\end{document}
